\chapter{Computers: Five-Loop, Fifth-Order Systems (5-OS)}\label{lect6a}
\localtoc


Let us recall that Alan Turing, defining the machine that bears his name, divided the tape into two parts when he introduced the concept of the Universal Turing Machine (UTM). One part for describing the automaton of the Turing machine $M(T)$ and another for the associated tape, $T$. In other words, the Turing Machine automaton can be viewed as operating with two memories, one for programs and one for data. At the same time, mathematicians also developed the concept of multi-tape TM, which involves two or more tapes with the same number of read heads, which allows accessing multiple symbols stored in different memories in the same execution cycle.

{\bf A little history}: on 29 March 1944 John von Neumann, who was working on the Manhattan Project, run one of the first programs  on the {\bf Harvard Mark I} computer. The architect of this computer, Howard Aiken, decided for his computer an organization with two memories, one for programs and another for data. Thus, the emergence of the {\bf {\em Harvard abstract model}} for computers inspired, like the non-Neumann abstract model, from the mathematical concept of the Turing Machine is known in history. 

\section{Interpreting vs. Executing}

Let us to go back to the UTM concept. At this theoretical level, the processing is done as a "collaboration" between the "program" from the tape and the finite automaton with its finite number of states. Each change on the "data" part of the UTM's tape is performed in a number of steps involving (1) readings form the "program" part of tape and (2) state transitions of the UTM's finite automaton. The complexity of the computational process is given by
\begin{enumerate}
    \item the size of "program", a flexible structure because it can be redesigned as many time as it is needed
    \item the size of the description of the UTM's finite automaton, a fixed structure because because it is frozen in a physical implementation.
\end{enumerate}
From the designing point of view, we care about the size of the automaton, more precisely, about the complexity of the associated combinational circuit. Any hardware complexity deserves to be reduced due to the effort involved in {\em specification, design, verification, validation} and {\em testing}. Consequently, minimizing the complexity of the control provided, at the origin of the UTM, by the finite automaton is welcome. Reducing the number of states of the finite automaton does not entail a reduction in complexity. This was proven by Claude Shannon's approach \cite{Shannon56}, which reduced the number of states to only 2, but with the consequence of increasing the size of the binary circuit of the loop that calculates the next state. Only the definition of a finite automaton with zero states \cite{Stefan06}, therefore the elimination of the finite automaton, justifies, only belatedly, the elimination of the complexity introduced by the control.

Consequently, we distinguish two ways of acting on an instruction read from program memory:
\begin{itemize}
    \item the instruction is {\bf {\em interpreted}} by executing a sequence of micro-instructions generated by a control automaton (see \ref{contrAut} in this book)
    \item the instruction is {\bf {\em executed}} in a single clock cycle using the instruction fields directly applied to the processor components. 
\end{itemize}
It is obvious that, for an executive structure, we consider the suggestion offered by a UTM with two tapes and a processing system originating in a UTM with zero internal states.


\section{Harvard Computer Version}

The Harvard abstract model assumes a fifth loop. The processor together with the program memory forms a 4-OS, and by closing another loop through the data memory, a 5-OS is obtained (see Figure \ref{fig45}).

\placedrawing[H]{figures/13_01_harvard_model.lp}{The Harvard abstract model  of computer.}{fig45}

We will continue to focus, for obvious reasons, on RISC-type processors (see \ref{toyrisc} in this book) that allow the implementation of the Harvard abstract model.

\placedrawing[H]{figures/13_02_computer_organization.lp}{The organization of a Harvard version of a computing system.}{compSys}

The main difference compared with the von Neumann abstract version is due to the dual connection of the processor to the memory system. Tho independent connections, one for data and another for programs, allow the processor to access in the same cycle an instruction and, if needed, a data word to be processed in the RALU subsystem. In this way the memory hierarchy deepens with severe consequences for memory access speed. A miss in the L1 cache is resolved much faster because a much faster memory is used than the main one. The same effect is produced by a third level of cache if used.

\subsection{The Reduction Instruction Set Computer (RISC) Approach}

When computer scientists become armed with huge amounts of performance data, they were in the position to demonstrate that the simple design of a RISC system was able to easily outperform even the most powerful classic CPU designs, while at the same time producing machine code that was only marginally larger than the heavily optimized CISC (Complex Instruction Set Computer) instructions.

A reduction instruction set maintains only the simplest and  the most frequently used instructions from a large instruction set, provided that all rejected instructions can be implemented with the sequences of the remaining ones.

\subsection{Segregation Between the Simple Hardware and the Complex Program}

The Harvard version of the computer allows for a clear segregation of the simple hardware structure from the complex program structure. Hardware is build mainly from simple circuits. See the hardware components of the toyRISC (see \ref{toyrisc} in this book) processor where only few gates in the decoder section (see \ref{interrupt}) can be considered complex circuits, i.e., with the size in the same range with the description size.


\section{Cache Memories}

There are three  levels of cache memory in the currently designed processors: Level 1 (L1), Level 2 (L2), and Level 3 (L3).

A cache memory is a small, volatile type of computer memory that provides high-speed  access to data and programs stored in the virtual memory space.  

The first level of cache, L1, is the smallest but fastest cache memory. It's connected directly with the processor on the same silicon die, and typically ranges from 8KB to 128KB\footnote{This quantitative information evolves very rapidly over time due to technological advances, which is why at the time of publication of this volume, it is certainly outdated.}.  L1  provides the fastest speed because it operates at the same speed as the processor. In a Harvard version of a computer L1 is divided into two parts: one is data cache, L1d, and another is instructions cache, L1i.

L2 is larger than L1 and slightly slower, but still faster than main memory. It can be located also on the processor. It's typically used to store data that's not as frequently accessed as the data in L1. The size of L2 cache memory varies  ranging from 1MB to 8MB.

The third level, L3, is  larger and slower than L2, but it's still faster than main memory. L3 cache memory is designed to improve the speed of L1 and L2. It can be found on the motherboard or, increasingly, inside the processor. Its size  ranges from 8MB to 64MB or more.

In summary, the different levels of cache memory serve to bridge the gap between the high speed of processor  and the main memory's lower speed. They store, according to the {\em locality principle} (see \ref{locprinc} in this book), frequently used data and instructions, which helps to speed up data access and overall system performance. The levels, according to the {\em memory hierarchy} (see \ref{memhier} in this book), differ in size, speed, and location, with each level being larger, slower, and generally further from the CPU than the previous one.

\section{A Little History: Abstract Models for Computers}

The following formulations have become established for abstract computer models: von Neumann architecture and Harvard architecture. First, let's see what the established descriptions are.

\paragraph{Harvard Architecture} was implicitly introduced in the design of the Harvard Mark I, a relay-based computer developed at Harvard University in 1944. The Harvard architecture was not formally named as such at the time but referred to a system where separate memory units were used for program instructions and data. This concept was largely attributed to the work of Howard Aiken, who led the development of the Harvard Mark I. The architecture became formally recognized as a distinct model in computing later, particularly in the 1970s-1980s, especially in contexts involving RISC processors.

\paragraph{Von Neumann Architecture} gained common usage in the 1960s-1970s and was introduced by John von Neumann in the 1945 report titled "First Draft of a Report on the EDVAC" \cite{vonNeumann1945} \cite{vonNeumann45}. This document outlined the fundamental principles of a stored-program computer, where program instructions and data share the same memory. The report was written as part of a collaborative effort during the development of the EDVAC (Electronic Discrete Variable Automatic Computer) at the Moore School of Electrical Engineering, University of Pennsylvania. Although von Neumann's name is most associated with the architecture, the theoretical contributions of Alan Turing, and the contributions of J. Presper Eckert, John Mauchly, and others on the the ENIAC and EDVAC projects were also significant.

\bigskip

The two models cannot be considered architectures in the sense in which the concept of architecture is commonly used in computer science and technology. We believe that a correct name, instead of architecture, is {\bf {\em abstract model}}. An abstract model establishes the connection between a mathematical model and a concrete realization, a physical structure. So: the {\bf Harvard abstract model} and the {\bf von Neumann abstract model}.

\subsection{Facts about Harvard Mark I Computer}

The Harvard Mark I, or IBM Automatic Sequence Controlled Calculator (ASCC), was a general-purpose electromechanical computers whose original concept was presented to IBM by Howard Aiken in November 1937. He developed the concept during his master's and doctoral studies at Harvard when he was faced with solving differential equations that could only be solved numerically. The project was approved to be founded by the company chairman Thomas Watson Sr. in February 1939. The ASCC was developed and built by IBM at their Endicott plant, in the state of New York, and shipped to Harvard in February 1944.

We can conclude that it is a computer made by IBM based on the project of Howard Aiken, a doctoral student at Harvard University. {\em Aiken -- IBM -- Harvard} is the chain linked to the Mark I computer that draws our attention in the attempt to use a name for the abstract model in question. We cannot help but notice that "Harvard" appears at the end of the aforementioned series.


\subsection{Facts about "First Draft"}

At the Moore School of Electrical Engineering at the University of Pennsylvania, Herman Goldstine, John Mauchly, J. Presper Eckert and Arthur Burks began to study the development of the EDVAC, a computing machine thought as the follower of the ENIAC. A serendipitous encounter of Goldstine with von Neumann led to the "First Draft" memo hand written by von Neumenn and sent by mail to Goldstine, who typed in 24 copies and distributed it with John von Neumann as the sole author.  

It should be noted that the material was not considered by von Neumann to be finalized. The title and author were given by Goldstine. The text does not contain any bibliographical references, which could have been the author's unfinished intention. Let's credit the author with good intentions, as these references were perhaps intended to be inserted into the text. So, {\em we cannot accuse the author} of omitting certain contributions, because no reference was actually made to the origin of the ideas contained in the report. But, at the same time, we cannot avoid mentioning the contributions of von Neumann's direct collaborators, to which must be added the theoretical contribution of Alan Turing, whom von Neumann knew and appreciated because he proposed him to be an assistant in one of his courses.

\subsection{Renamings}

We believe that the current use of the terms "Harvard architecture" and "von Neuman architecture" is inappropriate for several reasons.

First, the concept of architecture (see \ref{archOrg} in this book) has a definition, imposed in the 1960s \cite{Amdahl64}, with a well-defined meaning related to the functional interface between hardware and software. For this reason, it is much more appropriate to use the term "abstract model" instead of "architecture".

Second, the association of the abstract model with the name of a famous university or a renowned mathematician is not appropriate according to what has been described above. It can always be proven that other names are more appropriate.

We believe that we will be able to use more appropriate names if we also take into consideration aspects related to the internal structure of the processors involved in structures with a single external memory as opposed to those with two external memories. The distinction made in \ref{archOrg} between {\em interpreting processor} and {\em executing processor} is significant for the difference between the two types of abstract models. The interpreting processor (CISC) is controlled by a complex automaton that allows and requires the expansion of each instruction read from memory into a {\em micro-program}, while the executive processor (RISC) can use a single execution cycle for each instruction because it can simultaneously access two memories, the program and the data.

From the perspective of the above, we allow ourselves to propose the substitution of the phrase "Harvard architecture" with that of "{\em executive abstract model}", and the phrase "von Neumann architecture" with that of "{\em interpretive abstract model}". We thus obtain names that reflect more significant structural aspects than those related only to the organization of memory external to the processor.